Cine 심장 자기공명영상(cine CMR)은 한 심장주기를 여러 시간 프레임으로 나누어
촬영한 3차원 영상의 시계열이다. 이를 분리된 3차원 표면으로 변환하면 의료진이
원 영상과 움직이는 심장 모델을 나란히 보고, 회전시키거나 시간에 따른 형태 변화를
관찰할 수 있다. 그러나 각 프레임마다 2D Gaussian Splatting(2DGS)을 처음부터
최적화하면 전처리시간과 저장량이 프레임 수에 비례해 증가하고, primitive identity가
프레임마다 달라져 재생 시 흔들림이 발생한다.

본 논문은 CV-Dyn2DGS를 제안한다. 첫 프레임의 Chan--Vese 표면 위에 법선방향 두께가
없는 2차원 Gaussian 원반(surfel)을 배치하여 정준(canonical) 2DGS를 구축하고, 이후
프레임에서는 이전 level-set을 seed로 현재 표면을 구한 뒤 기존 surfel의 중심을
새 표면의 법선 방향으로 정사영하고 접평면 프레임을 최소회전으로 전송한다. 표면
이동만으로 설명되지 않는 밝기 변화는 작은 프레임별 외형 잔차로 보정한다. 저장되는
것은 정준 2DGS 하나, 프레임별 표면, 그리고 작은 잔차뿐이며, dense deformation
field나 학습된 motion field는 저장하지 않는다.

이론적 기여는 네 가지이다. 첫째, 비등방 복셀 간격 $\spacing=(h_x,h_y,h_z)$을 명시적으로
반영한 간격 인지 이산화를 정립한다. 단기축 cine CMR은 $h_z \gg h_x \approx h_y$이므로
간격을 무시한 스텐실은 곡률 기반 정규화를 슬라이스 방향으로 편향시킨다. 둘째, 인접
프레임 해의 근접성을 이용한 좁은 띠 웜스타트가 수렴을 보존하며 반복 횟수를 줄임을
국소 Polyak--{\L}ojasiewicz 조건 아래 보인다. 셋째, 일차 Taylor 제약에 대한 최소
노름 논법으로 최소법선 보정을 유도하고, 그 반복이 잔차를 이차 속도로 감소시킴을
보인다. 넷째, 접평면 프레임, 원근 정합 래스터화, 외형 잔차 최소제곱, 중간 시점
보간의 네 축에 대해 명시적 오차 한계를 유도하고 각 한계를 검증 지표에 대응시킨다.

본 논문은 위 이론을 PyTorch로 구현한 시스템을 함께 제시한다. 구현은 간격 인지
연산자, 좁은 띠 Chan--Vese, 정준 surfel 초기화, 법선 정사영, 접선 전송, 원근 정합
타일 래스터화, 희소 렌더링 연산자를 이용한 정칙화 최소제곱, 저계수 잔차 압축,
그리고 mesh 및 thin-3DGS 비교 기준선을 포함한다.

\medskip
\noindent\textbf{실험적 주장의 현재 상태.}
본 논문의 효율 및 품질 주장 RQ1--RQ6은 \emph{가설이며 아직 확인되지 않았다}.
구현은 완성되었으나 실행된 적이 없다(제 \ref{sec:impl-status}절). 따라서 제 8장의
모든 표는 자리표시자이며, 어떤 성능 수치도 제시하지 않는다. 다만 PyTorch를
필요로 하지 않는 이산 논리 커널 41개 항목은 독립 참조와 대조하여 검증되었다
(제 \ref{sec:kernel-verification}절).
